\documentclass[preview]{standalone}

\usepackage{amsmath}
\usepackage{amssymb}
\usepackage{stellar}
\usepackage{definitions}

\begin{document}

\id{la-ricerca-in-psicologia}
\genpage

\section{La ricerca in psicologia}

\begin{snippetdefinition}{psicologia-ingenua-definition}{Psicologia ingenua}
    La \emph{psicologia ingenua}, o psicologia del senso comune, spiega il
    comportamento attraverso teorie fondate sull'esperienza personale e non sul
    controllo scientifico.
\end{snippetdefinition}

\begin{snippet}{processo-ricerca-psicologia}
    Il processo di ricerca in psicologia serve a determinare quali supposizioni
    sul comportamento siano corrette, verificarle e approfondirle. Le domande
    dei ricercatori nascono da osservazioni, credenze, aspetti culturali e
    idee formulate in precedenza.
\end{snippet}

\includesnpt[width=75\%,src=/snippet/static-psychology/psychology-research-process.jpg]{centered-img}

\begin{snippetdefinition}{teoria-psicologia-definition}{Teoria}
    Una \emph{teoria} è un insieme organizzato di proposizioni che spiegano un
    fenomeno o un insieme di fenomeni. In psicologia, una teoria fornisce il
    contesto entro cui formulare le domande di ricerca.
\end{snippetdefinition}

\begin{snippetdefinition}{ipotesi-psicologia-definition}{Ipotesi}
    Un'\emph{ipotesi} è un'affermazione provvisoria e verificabile sulla
    relazione tra cause e conseguenze di un fenomeno. Per essere scientifica,
    deve essere formulata in modo tale da poter essere falsificata.
\end{snippetdefinition}

\begin{snippetexample}{ipotesi-zuccheri-example}{Ipotesi sugli zuccheri}
    L'affermazione ``se i bambini consumano molti zuccheri, allora è probabile
    che manifestino comportamenti iperattivi'' è un'ipotesi verificabile. Per
    controllarla, si possono far consumare zuccheri a un gruppo di bambini e
    osservare eventuali cambiamenti nei livelli di attivazione. Se non si
    trova alcun effetto, l'ipotesi può essere respinta.
\end{snippetexample}

\begin{snippetdefinition}{metodo-scientifico-psicologia-definition}{Metodo scientifico}
    Il \emph{metodo scientifico} è un insieme generale di procedure per la
    raccolta e l'interpretazione dei dati, finalizzato a limitare le fonti di
    errore e a trarre conclusioni verificabili.
\end{snippetdefinition}

\begin{snippetdefinition}{metodo-ricerca-psicologia-definition}{Metodo di ricerca}
    Un \emph{metodo di ricerca} è una procedura scientifica che consente di
    raccogliere e mettere in rapporto dati empirici, in modo da poterli
    interpretare rispetto a un'ipotesi.
\end{snippetdefinition}

\begin{snippetdefinition}{dati-grezzi-definition}{Dati grezzi}
    I \emph{dati grezzi} sono informazioni raccolte attraverso procedure
    prestabilite, cioè protocolli sperimentali, e possono essere usati come
    prove a sostegno o contro le ipotesi.
\end{snippetdefinition}

\begin{snippet}{tipi-dati-psicologia}
    In psicologia si distinguono due tipi principali di dati:
    \begin{itemize}
        \item i \emph{dati comportamentali}, collegati alla rilevazione del
        comportamento umano o animale;
        \item i \emph{dati self-report}, cioè descrizioni che le persone
        forniscono riguardo a ciò che accade nella loro mente.
    \end{itemize}
\end{snippet}

\begin{snippet}{analisi-pubblicazione-psicologia}
    Una volta raccolti i dati, i ricercatori li analizzano e accettano o
    falsificano l'ipotesi. L'analisi statistica usa statistiche descrittive,
    per sintetizzare i dati, e statistiche inferenziali, per formulare giudizi
    e previsioni. Se i risultati hanno valore scientifico, vengono pubblicati
    in modo che la comunità possa verificarli.
\end{snippet}

\end{document}
